%\documentclass[10pt]{article}
%\usepackage[latin1]{inputenc}
%\usepackage[frenchb]{babel}
%\usepackage[T1]{fontenc}
%\usepackage{listings}
%\usepackage[colorlinks=true,urlcolor=blue]{hyperref}
%\usepackage[vmargin=2cm,hmargin=2cm]{geometry}
%\usepackage[svgnames,table]{xcolor}
%\usepackage{kpfonts}
%\usepackage{cellspace}
%\usepackage{aurical}
%\usepackage{graphicx}
%\RequirePackage{titlesec}
%\usepackage[electrum]{pas-keyboard}
% <---------- Param�tres globaux -------------------->
%\author{St�phane PASQUET \\ \href{http://www.mathweb.fr}{http://www.mathweb.fr} \\ \href{mailto:contact@mathweb.fr}{contact@mathweb.fr}}
%\date{\today}
%\setlength{\parindent}{0cm}
%\definecolor{touchColor}{cmyk}{0,0,.05,0}
%\pagecolor{touchColor}
%\title{Documentation\\ \texttt{pas-keyboard.sty V1.02}}

\documentclass[a4paper,french]{article}
\usepackage[latin1]{inputenc}
\usepackage[T1]{fontenc}
\usepackage{babel} 
\usepackage{etex}
\usepackage{fourier}
\usepackage[table]{xcolor}
\usepackage[colorlinks=true,urlcolor=blue]{hyperref}
\usepackage{pas-keyboard}
\usepackage{titlesec}
\usepackage{tcolorbox}
	\tcbuselibrary{skins}
	\tcbuselibrary{theorems}
	\tcbuselibrary{breakable}

% --- Propre � cette doc

\graphicspath{{images/}}
\usepackage{longtable}
\usepackage{tabularx}
\usepackage{cellspace}
\setlength{\cellspacetoplimit}{4pt}
\setlength{\cellspacebottomlimit}{4pt}

% ----------------------

\setlength{\parindent}{0pt}

\tcbuselibrary{listings}
\usetikzlibrary{decorations.pathmorphing}

% Couleurs utilisées dans la documentation

\definecolor{codeTitleFont}{cmyk}{0.04,0,0.03,0.16}
\definecolor{codeTitleBackLeft}{cmyk}{0.08,0,0.06,0.76}
\definecolor{codeTitleBackRight}{cmyk}{0.07,0,0.05,0.42}
\definecolor{listingTitleFont}{cmyk}{0,0.31,0.91,0.38}
\definecolor{listingTitleBackLeft}{cmyk}{0,0.05,0.64,0}
\definecolor{listingTitleBackRight}{cmyk}{0,0.03,0.31,0.02}


% Code LaTeX

\tcbset{codeTEX/.style={
	sharp corners=all,
	before skip=1em,
	after skip=1em,
	enhanced,
	frame style={	
			left color=codeTitleBackLeft,
			right color=codeTitleBackRight},
	interior style={
		top color=codeTitleBackLeft!50,
		bottom color=codeTitleBackRight!20},
	boxrule=0.7pt,
	fonttitle={\sffamily\bfseries\color{codeTitleFont}},
	colback=codeTitleFont,
	listing only,
	left=6mm,
	listing options={
		basicstyle=\ttfamily\fontsize{7}{9}\selectfont,
		keywordstyle=\color{blue},
		numbers=left,
		language=TeX,
		breaklines=true,
		morekeywords={definecolor,tcbset,begin, ssfamily,Huge,newenvironment,newcommand,bfseries,color, usepackage,tcblower,ttfamily,documentclass,newtitlesubject, newtitlecorrection,titlesubject,titlecorrection,newpageforcorrection ,newpageforsubject},
		numberstyle=\tiny\color{red!75!black}},
	breakable
	}
}

% Listing exemples

\tcbset{listing/.style={
	sharp corners=all,
	before skip=1em,
	after skip=1em,
	enhanced,
	frame style={	
			left color=listingTitleBackLeft,
			right color=listingTitleBackRight},
	boxrule=0.7pt,
	fonttitle={\sffamily\bfseries\color{listingTitleFont}},
	colback=listingTitleBackRight,
	breakable,
	listing options={
		basicstyle=\ttfamily\fontsize{7}{9}\selectfont,
		keywordstyle=\color{listingTitleFont},
		numbers=left,
		language=TeX,
		breaklines=true,
		numbersep=5pt,
		morekeywords={ifelse,begin,definecolor,tcbset},
		numberstyle=\tiny\color{red!75!black}},
	},
	interior style={
		draw=listingTitleBackLeft,
		top color=listingTitleBackLeft!50,
		bottom color=listingTitleBackRight!20},
	  segmentation style={
		draw=listingTitleFont,
		solid,
		decorate,
		decoration={random steps,segment length=2mm}
	}
}

% Titre de la documentation

\tcbset{head/.style={
	enhanced,
	hbox,
	tikznode,
	left=8mm,
	right=8mm,
	boxrule=0.4pt,
  colback=white,
  colframe=gray,
  drop lifted shadow=black!50!yellow,
  before=\par\vspace*{5mm},
  after=\par\bigskip,
  interior style={
		draw=white,
		top color=white,
		bottom color=white}
	}
}

% TOC

\tcbset{toc/.style={
	breakable,
	enhanced jigsaw,
	title={\color{red!50!black}Sommaire},
	fonttitle=\bfseries\Large,
  colback=yellow!10!white,
  colframe=red!50!black,
  before=\par\bigskip\noindent,
  interior style={
  	fill overzoom image=goldshade.png,
  	fill image opacity=0.25},
  colbacktitle=yellow!20,
  enlargepage flexible=\baselineskip,
  pad at break*=3mm,
  attach boxed title to top center={
  	yshift=-0.25mm-\tcboxedtitleheight/2,
  	yshifttext=2mm-\tcboxedtitleheight/2},
  boxed title style={
  	enhanced,
  	boxrule=0.5mm,
    frame code={ 
    \path[tcb fill frame] ([xshift=-4mm]frame.west) -- (frame.north west)
    -- (frame.north east) -- ([xshift=4mm]frame.east)
    -- (frame.south east) -- (frame.south west) -- cycle; },
    interior code={ 
    	\path[tcb fill interior] ([xshift=-2mm]interior.west)
    -- (interior.north west) -- (interior.north east)
    -- ([xshift=2mm]interior.east) -- (interior.south east) -- (interior.south west)
    -- cycle;}  },
  drop fuzzy shadow
	}
}

% Historique de l'extension

\tcbset{histo/.style={
	enhanced,
	breakable,
	sidebyside,
	lefthand width=1.8cm
	}
}
\makeatletter

\def\@dottedtocline#1#2#3#4#5{%
  \ifnum #1>\c@tocdepth \else
    \vskip \z@ \@plus.2\p@
    {\leftskip #2\relax \rightskip \@tocrmarg \parfillskip -\rightskip
     \parindent #2\relax\@afterindenttrue
     \interlinepenalty\@M
     \leavevmode
     \@tempdima #3\relax
     \advance\leftskip \@tempdima \null\nobreak\hskip -\leftskip
     {#4}\nobreak
     \leaders\hbox{$\m@th
        \mkern \@dotsep mu\hbox{.}\mkern \@dotsep
        mu$}\hfill
     \nobreak
     \hb@xt@\@pnumwidth{\hfil\helvbx #5}%
     \par}%
  \fi}

\renewcommand*\l@section
{%
\helvbx\color{red!50!black}\bfseries
\def\@linkcolor{red!50!black}\@dottedtocline{1}{1.5em}{1.5em}
}

\renewcommand*\l@subsection
{%
\helvbx\color{green!50!black}
\def\@linkcolor{green!50!black}
\@dottedtocline{1}{2.3em}{2.6em}
}

\renewcommand*\l@subsubsection
{%
\helvbx\color{orange!80!black}
\def\@linkcolor{orange!80!black}
\@dottedtocline{1}{3em}{3.3em}
}

\def\contentsline#1#2#3#4{%
  \ifx\\#4\\%
    \csname l@#1\endcsname{#2}{#3}%
  \else
      \csname l@#1\endcsname{\hyper@linkstart{link}{#4}{#2}\hyper@linkend}{%
        \hyper@linkstart{link}{#4}{#3}\hyper@linkend
      }%
  \fi
}

% --------------------
% TITRES DES SECTIONS 
% --------------------

\titleformat{\section}[block]
{\helvbx\Large\color{red!50!black}}
{\fcolorbox{red!50!black}{red!50!black}{\textcolor{white}{\bfseries\thesection}}}
{1em}
{\helvbx}

\titleformat{\subsection}[block]
{\helvbx\large\color{green!50!black}}
{\thesubsection}
{1em}
{\helvbx}

\titleformat{\subsubsection}[block]
{\helvbx\large\color{orange!50!black}}
{\thesubsubsection}
{1em}
{\helvbx}

\makeatother

\newcommand{\helvbx}{\usefont{T1}{phv}{m}{n}}

\begin{document}

\begin{center}
\begin{tcolorbox}[head]
{\bfseries\LARGE Documentation \texttt{pas-keyboard} }\\[3mm]
{\large Version 1.02 -- \today}
\end{tcolorbox}

{\large 
\href{http://www.mathweb.fr/contact.html}{St�phane Pasquet}}
\end{center}

\begin{tcolorbox}[toc]
\makeatletter
\@starttoc{toc}
\makeatother
\end{tcolorbox}

\section{Introduction et installation}

L'extension \texttt{pas-keyboard} est con�ue pour afficher des touches d'ordinateur.\\
Bien s�r, il y a l'extension \texttt{keystroke}, mais le rendu me ne satisfaisait pas. C'est la raison pour laquelle j'ai eu envie de cr�er cette extension.

\medskip

Cette extension charge automatiquement les extensions suivantes :

\medskip

\begin{quote}
tikz (avec les librairies : calc, arrows)
\end{quote}

\medskip

Sous Ubuntu, on pourra d�compresser \texttt{pas-crosswords.zip} dans le r�pertoire :

\begin{verbatim}
./texlive/texmf-local/tex/latex/
\end{verbatim}

de sorte � avoir :


\begin{verbatim}
./texlive/texmf-local/tex/latex/pas-keyboard/latex/pas-keyboard.sty
./texlive/texmf-local/tex/latex/pas-keyboard/doc/pas-keyboard.tex
./texlive/texmf-local/tex/latex/pas-keyboard/doc/pas-keyboard.pdf
./texlive/texmf-local/tex/latex/pas-keyboard/doc/doc.codes.tex
./texlive/texmf-local/tex/latex/pas-keyboard/doc/doc.styles.tex
\end{verbatim}

\medskip

Apr�s installation, n'oubliez pas de taper la commande \texttt{texhash} dans le terminal pour mettre � jour la base de donn�es des extensions.

\bigskip 

Avec Miktex (sous Windows) ou macTex (sous Mac OS), j'imagine que l'arborescence ressemble � ce qui est �crit pr�c�demment.

\section{Appel de l'extension}

Afin d'imiter le mieux les touches d'un ordinateur, j'ai fait appel � la fonte \texttt{electrum} disponible sur la page :
\begin{center}
\url{http://www.ctan.org/tex-archive/fonts/electrumadf/}. 
\end{center}

Cependant, conscient des difficult�s que pose la mise en place d'une telle fonte, j'ai d�cid� de faire appel � celle-ci uniquement sur l'appel d'une option. Ainsi, pour faire appel � cette police de caract�res, l'extension devra �tre appel�e avec l'option suivante :

\medskip

\begin{tcblisting}{codeTEX}
\usepackage[electrum]{pas-keyboard}
\end{tcblisting}

\medskip

Sans cette option, les caract�res des touches seront mis en helvetica.

\section{Les commandes}

Voici une liste des diff�rentes commandes :

\medskip\arrayrulecolor{brown}

\begin{center}
\begin{longtable}{|>{\centering\arraybackslash}Sc|>{\centering\arraybackslash}Sc|}
\hline\rowcolor{brown}\color{white}
Commandes & \color{white}R�sultats\\
\hline
\texttt{\textbackslash key$\{$M$\}$}& \key{M}\\
\hline
\texttt{\textbackslash keyCTRL}& \keyCTRL\\
\hline
\texttt{\textbackslash keyUpArrow}& \keyUpArrow\\
\hline
\texttt{\textbackslash keyDownArrow}& \keyDownArrow\\
\hline
\texttt{\textbackslash keyLeftArrow}& \keyLeftArrow\\
\hline
\texttt{\textbackslash keyRightArrow}& \keyRightArrow\\
\hline
\texttt{\textbackslash keyShift}& \keyShift\\
\hline
\texttt{\textbackslash keySHIFT}& \keySHIFT\\
\hline
\texttt{\textbackslash keyBack}& \keyBack\\
\hline
\texttt{\textbackslash keyBACK}& \keyBACK\\
\hline
\texttt{\textbackslash keyEnter}& \keyEnter\\
\hline
\texttt{\textbackslash keyTab}& \keyTab\\
\hline
\texttt{\textbackslash keyCapsOff}& \keyCapsOff\\
\hline
\texttt{\textbackslash keyCapsOn}& \keyCapsOn\\
\hline
\texttt{\textbackslash keySupInf}& \keySupInf\\
\hline
\texttt{\textbackslash keySquare}& \keySquare\\
\hline
\texttt{\textbackslash keyEsc}& \keyEsc\\
\hline
\texttt{\textbackslash keyHome}& \keyHome\\
\hline
\texttt{\textbackslash keyHOME}& \keyHOME\\
\hline
\texttt{\textbackslash keyPgUp}& \keyPgUp\\
\hline
\texttt{\textbackslash keyPgDown}& \keyPgDown\\
\hline
\end{longtable}
\end{center}

\bigskip

Bien entendu, vous pouvez mettre n'importe quelle lettre dans la commande \texttt{\textbackslash key} :

\medskip

\begin{tcblisting}{listing}
\key{A}\key{B}\key{C}
\end{tcblisting}

\section{Les param�tres}

Vous avez la possibilit� de changer les couleurs :

\begin{center}
\begin{tabularx}{\linewidth}{|>{\centering\arraybackslash}p{2cm}|>{\centering\arraybackslash}X|p{6.5cm}|}
\hline\rowcolor{brown}
\color{white}Nom & \color{white}Correspondance & \color{white}Valeur par d�faut\\
\hline
\texttt{touchColor} & Fond des touches & \textbackslash\texttt{definecolor$\{$touchColor$\}\{$cmyk$\}\{$0,0,0,0$\}$}\\
\hline
\texttt{textColor} & Texte des touches & \textbackslash\texttt{definecolor$\{$textColor$\}\{$cmyk$\}\{$1,1,1,1$\}$}\\
\hline
\texttt{drawColor} & Contours des touches & \textbackslash\texttt{definecolor$\{$drawColor$\}\{$cmyk$\}\{$1,1,1,1$\}$}\\
\hline
\texttt{shadowColor} & Ombre des touches & \textbackslash\texttt{definecolor$\{$shadowColor$\}\{$cmyk$\}$ $\{$0.30,0.20,0.20,0$\}$}\\
\hline
\texttt{lightON} & Couleur de la \og diode \fg{} de la touche CAPS LOCK quand elle est active & \textbackslash\texttt{definecolor$\{$lightON$\}\{$cmyk$\}$ $\{$0.81,0.11,1,0.02$\}$}\\
\hline
\texttt{lightOFF} & Couleur de la \og diode \fg{} de la touche CAPS LOCK quand elle est inactive & \textbackslash\texttt{definecolor$\{$lightOFF$\}$ $\{$cmyk$\}\{$0.55,0.42,0.42,0.09$\}$}\\
\hline
\end{tabularx}
\end{center}

\medskip

Vous pouvez donc changer la couleur de votre choix en la red�finissant avec la commande \textbackslash\texttt{definecolor}.

\section{Suppl�ments}

Vous pouvez ins�rer ce que vous souhaitez dans les touches, y compris des dessins comme le montre le script suivant :

\medskip

\begin{tcblisting}{listing}
\newcommand{\dessin}{\tikz{\clip (0,-0.1) 
rectangle +(.4,.2);\draw (0,0) .. controls 
(.2,.3) and (.3,-.3) .. (.4,.2);}}
\key{\dessin}
\end{tcblisting}

\medskip

Mais vous n'�tes pas oblig�s d'ins�rer du code TiKZ ; un code PST ou m�me une image est possible comme le montre le script suivant :

\medskip

\begin{tcblisting}{listing}
\key{\includegraphics[scale=0.5]{icone.png}}
\end{tcblisting}

\medskip

Je vous rappelle que la commande \textbackslash\texttt{includegraphics} s'emploie en ayant appel� l'extension \texttt{graphicx}.

\section{Mises � jour}

\begin{tcolorbox}[histo]
2013/10/03
~
\tcblower
Modification de la commande \texttt{\textbackslash key} afin de centrer verticalement les touches sur la ligne courante.
\end{tcolorbox}
\end{document}